\documentclass[12pt]{article}
\usepackage[T1]{fontenc}
\usepackage[utf8]{inputenc}
\usepackage{lmodern}
\usepackage{textcomp}
\usepackage{enumitem}
\usepackage{geometry}
\usepackage{graphicx}
\usepackage{amsmath, amssymb}
\geometry{margin=1in}
\begin{document}
\begin{center}
Sociology - Undergraduate Level Assessment 23 \
Total Marks: 40 \
Answer all questions.
\end{center}

\section*{Multiple Choice (10 marks)}
\begin{enumerate}[label=\arabic*.]
\item The alternative form of 'network organization' means that:
\begin{enumerate}[label=(\alph*)]
    \item work can be subcontracted out to independent suppliers and retailers
    \item business transactions occur only through electronic communication
    \item the Japanese model is applied, through lateral networks of flexible roles
    \item activities are redistributed equally between men and women
\end{enumerate}
\item In the nineteenth century, homosexuality was understood as:
\begin{enumerate}[label=(\alph*)]
    \item a positive identity in which gay people could take pride
    \item an absolute taboo, which meant that all homosexuals were isolated
    \item a subordinate form of masculinity that threatened 'compulsory heterosexuality'
    \item confirmation of the two-sex model
\end{enumerate}
\item The capitalist world economy is what Wallerstein (1974) would call a 'world system'. This term refers to:
\begin{enumerate}[label=(\alph*)]
    \item a means of transporting money between different areas of a country
    \item an empire with a bureaucratic administration but no political centre
    \item an awareness of risks and dangers that affect the environment as a whole
    \item a unit with a division of labour that extends across ethnic and cultural groups
\end{enumerate}
\item Sreberny-Mohammadi (1996) argues that national cultures can resist American cultural domination of the media by:
\begin{enumerate}[label=(\alph*)]
    \item domesticating its content, including more 'home-produced' programmes
    \item controlling the distribution of imported products by banning satellite dishes
    \item creating 'reverse flows' of their own programmes back to imperial societies
    \item all of the above
\end{enumerate}
\item Walt Disney, Sony and Time Warner are examples of:
\begin{enumerate}[label=(\alph*)]
    \item transnational corporations
    \item multi-media empires
    \item ownership concentrated within one medium
    \item government-owned companies
\end{enumerate}
\end{enumerate}

\section*{True / False (10 marks)}
\textit{Answer True or False}
\begin{enumerate}[label=\arabic*.]
\setcounter{enumi}{5}
\item True or False: Government 'white papers' are considered personal documents with restricted access due to their confidential nature.
\item True or False: The social construction of childhood can be traced back to compulsory education, emotional ties, and consumer goods for children.
\item True or False: Privatization led to employment being contingent on performance rather than on trust and commitment between employees and managers.
\item True or False: In stage 3 of the health transition, chronic degenerative diseases like cancer and heart disease are the primary causes of illness and death.
\item True or False: Dahrendorf, Rex, and Habermas primarily examined social solidarity and cohesion in their sociological theories.
\end{enumerate}

\section*{Long Form Response (20 marks)}
\textit{Answer each question with a clear and complete explanation.}
\begin{enumerate}[label=\arabic*.]
\setcounter{enumi}{10}
\item Discuss how differential association theory explains the influence of peer behavior on individual choices, using the analogy of friends jumping off a bridge as a starting point for your analysis.
\item Discuss how the functionalist perspective in sociology conceptualizes society as a system of interdependent and coordinated parts, and analyze its implications for social stability and change.
\end{enumerate}

\vfill
\noindent\textit{End of Paper}
\end{document}